\section{Discussion}
\label{sec:discussion}

The results support a self-test in two levels for the realistic circuit. The compliance decision
can be set to let through about \SI{1}{\percent} of the non-compliant circuits while rejecting
\SI{8}{\percent} of the compliant ones, or \SI{2}{\percent} when porous dry electrodes are
excluded. A limit test cannot approach this, because it asks whether the circuit looks like a
healthy one, which is neither necessary nor sufficient for compliance. The second level should
report an ambiguity group, not a component: the groups are known from the nominal circuit, they
hold for fault magnitudes the model has not seen, and in \cA\ a group names two to four parts
that set the same quantity. The comparison of labels has a direct practical reading: a detector
trained in the usual way of the diagnosis literature would take out of service more than a third
of the instruments that still meet the standard. The cost is low: three measurements and three
seconds give the decision, and the tones below the high-pass corner are not needed, which
matters because the self-test runs with the patient connected.

The electrode limits both levels. A healthy circuit on a high-impedance electrode measures a
reduced gain through the calibration source, which is what a gain fault looks like; on \cA\ this
accounts for eight of the nine invisible fault conditions, and false rejects fall from 7.8 to
\SI{1.9}{\percent} without porous electrodes. A
degraded contact can only be recognized against the normal contact of the same subject, which a
single self-test does not know. A self-test that kept a reference per patient, or that could
disconnect the electrodes from the calibration path, would remove most of the difficulty;
neither was studied.

The two circuits differ only in how the instrumentation amplifier is built, and that changes the
conclusions. The discrete amplifier adds seven resistors that all act on the gain and several
DRL faults that are non-compliant and invisible; evaluated only on it, the method would have been
judged unable to reach a low escape rate. This supports the argument of~\cite{chen2025} that the
circuit conditions the conclusions, with a different remedy: the ambiguity of the given circuit
is computed and the diagnosis reported at that level. The invisible DRL faults also show a limit
of the measurements: where a spare converter channel is available, a dc reading of the DRL output
is worth more than a larger common-mode tone.

\emph{Limitations.} The study is based on simulation, with behavioral amplifier models and ideal
self-test sources; transfer to hardware has not been tested. Faults are single and aging is not
modeled. The dry-electrode model rests on six subjects per material, and the gel model and the
spread between electrodes are assumed. The noise specification is computed from the rms value
and the input range from the operating point. The specifications are those of performance: some
hard faults leave the circuit compliant while removing a protection (a short across the resistor
that limits the current into the right-leg electrode), and a self-test aimed at safety would
have to label them separately.
